\section{总结与展望}

\subsection{工作总结}

本项目实现了一个面向生物医学数据科学的推理编码智能体系统，主要完成了以下工作：

\begin{enumerate}
    \item \textbf{ReAct 核心智能体}：实现了完整的 Thought--Action--Observation 闭环控制逻辑，包含持久化执行命名空间、主动验证机制和渐进式提示策略，在 BioDSBench 数据集上取得了良好的任务通过率。
    \item \textbf{反思与经验复用池}：设计了结构化的经验存储、LLM 驱动的反思引擎和基于分析类型的经验检索机制。消融实验表明，该机制在 qwen-turbo 模型上带来了 10\% 的通过率提升。
    \item \textbf{技能库与 AST 指纹检索}：提出了"Agent as Tool Maker"机制，通过 AST 调用链指纹提取和 LCS 相似度匹配实现代码级知识复用。技能库与经验池形成互补，联合覆盖率达 73.3\%，比单独使用任一机制高 6.6 个百分点；在共同通过的任务上，技能库的平均步骤数减少了 8.9\%。
    \item \textbf{多智能体协作}：构建了包含四类专门化智能体的协作系统，通过消息总线实现异步通信，支持顺序、并行、分层和自适应四种协作模式。
    \item \textbf{Web 可视化界面}：基于 Streamlit 开发了实时监控面板，支持任务配置、执行监控和结果导出。
\end{enumerate}

在 30 道 BioDSBench 任务的五组消融实验中，qwen3-max 完整配置达到 90.0\% 的通过率，qwen-turbo 在经验池辅助下达到 66.7\%，技能库（预热）同样达到 66.7\% 且平均步骤更少，验证了系统设计的有效性和两种知识复用机制的互补价值。

\subsection{局限性}

本项目存在以下局限性：

\begin{enumerate}
    \item \textbf{经验检索精度}：当前基于关键词匹配的检索方式较为粗糙，对于语义相似但表述不同的任务可能检索不到相关经验。引入向量嵌入的语义检索可以改善这一问题。
    \item \textbf{技能库冷启动}：技能库在初始阶段没有可用技能，需要预热或经历一定数量的任务积累后才能发挥作用。冷启动组（60.0\%）低于预热组（66.7\%）的结果证实了这一局限。
    \item \textbf{多智能体协作深度}：当前的多智能体系统主要在任务级别进行分工，尚未实现步骤级别的细粒度协作（如一个智能体写代码、另一个智能体审查代码）。
    \item \textbf{沙箱序列化开销}：安全沙箱通过 pickle 在主进程和子进程间传递命名空间，对于大型 DataFrame 等对象存在序列化/反序列化的性能开销。
\end{enumerate}

\subsection{未来工作}

\begin{enumerate}
    \item \textbf{经验池与技能库融合}：当前两种机制独立运行，未来可设计统一的知识管理框架，根据任务特征自适应选择注入经验（few-shot）还是技能（可执行函数），或两者联合注入。实验表明联合覆盖率可达 73.3\%，融合策略有望进一步提升。
    \item \textbf{语义经验检索}：使用文本嵌入模型（如 text-embedding-v3）对经验和技能进行向量化，通过余弦相似度实现更精准的语义检索，替代当前的关键词匹配。
    \item \textbf{在线学习}：在 benchmark 运行过程中实时更新经验池和技能库，使后续任务能够立即受益于前序任务的积累。
    \item \textbf{代码审查智能体}：在多智能体系统中引入专门的代码审查角色，在代码执行前进行静态分析和逻辑检查，减少运行时错误。
    \item \textbf{更大规模评估}：在完整的 108 道 BioDSBench 任务上进行评估，并与已发表的 baseline 方法进行对比。
\end{enumerate}
